%% bare_jrnl.tex
%% V1.4b
%% 2015/08/26
%% by Michael Shell
%% see http://www.michaelshell.org/
%% for current contact information.
%%
%% This is a skeleton file demonstrating the use of IEEEtran.cls
%% (requires IEEEtran.cls version 1.8b or later) with an IEEE
%% journal paper.
%%
%% Support sites:
%% http://www.michaelshell.org/tex/ieeetran/
%% http://www.ctan.org/pkg/ieeetran
%% and
%% http://www.ieee.org/

%%*************************************************************************
%% Legal Notice:
%% This code is offered as-is without any warranty either expressed or
%% implied; without even the implied warranty of MERCHANTABILITY or
%% FITNESS FOR A PARTICULAR PURPOSE! 
%% User assumes all risk.
%% In no event shall the IEEE or any contributor to this code be liable for
%% any damages or losses, including, but not limited to, incidental,
%% consequential, or any other damages, resulting from the use or misuse
%% of any information contained here.
%%
%% All comments are the opinions of their respective authors and are not
%% necessarily endorsed by the IEEE.
%%
%% This work is distributed under the LaTeX Project Public License (LPPL)
%% ( http://www.latex-project.org/ ) version 1.3, and may be freely used,
%% distributed and modified. A copy of the LPPL, version 1.3, is included
%% in the base LaTeX documentation of all distributions of LaTeX released
%% 2003/12/01 or later.
%% Retain all contribution notices and credits.
%% ** Modified files should be clearly indicated as such, including  **
%% ** renaming them and changing author support contact information. **
%%*************************************************************************


% *** Authors should verify (and, if needed, correct) their LaTeX system  ***
% *** with the testflow diagnostic prior to trusting their LaTeX platform ***
% *** with production work. The IEEE's font choices and paper sizes can   ***
% *** trigger bugs that do not appear when using other class files.       ***                          ***
% The testflow support page is at:
% http://www.michaelshell.org/tex/testflow/



\documentclass[journal]{IEEEtran}
%
% If IEEEtran.cls has not been installed into the LaTeX system files,
% manually specify the path to it like:
% \documentclass[journal]{../sty/IEEEtran}
\usepackage{url}
\usepackage{multirow}
\usepackage{amsmath}
\usepackage{booktabs}
\usepackage{amssymb}
\usepackage{float}
\usepackage{graphicx} 
\usepackage{adjustbox} 
\usepackage{caption}
\usepackage{subcaption}

\usepackage{xcolor}
\usepackage{comment}
\usepackage{makecell}
% Some very useful LaTeX packages include:
% (uncomment the ones you want to load)


% *** MISC UTILITY PACKAGES ***
%
%\usepackage{ifpdf}
% Heiko Oberdiek's ifpdf.sty is very useful if you need conditional
% compilation based on whether the output is pdf or dvi.
% usage:
% \ifpdf
%   % pdf code
% \else
%   % dvi code
% \fi
% The latest version of ifpdf.sty can be obtained from:
% http://www.ctan.org/pkg/ifpdf
% Also, note that IEEEtran.cls V1.7 and later provides a builtin
% \ifCLASSINFOpdf conditional that works the same way.
% When switching from latex to pdflatex and vice-versa, the compiler may
% have to be run twice to clear warning/error messages.






% *** CITATION PACKAGES ***
%
%\usepackage{cite}
% cite.sty was written by Donald Arseneau
% V1.6 and later of IEEEtran pre-defines the format of the cite.sty package
% \cite{} output to follow that of the IEEE. Loading the cite package will
% result in citation numbers being automatically sorted and properly
% "compressed/ranged". e.g., [1], [9], [2], [7], [5], [6] without using
% cite.sty will become [1], [2], [5]--[7], [9] using cite.sty. cite.sty's
% \cite will automatically add leading space, if needed. Use cite.sty's
% noadjust option (cite.sty V3.8 and later) if you want to turn this off
% such as if a citation ever needs to be enclosed in parenthesis.
% cite.sty is already installed on most LaTeX systems. Be sure and use
% version 5.0 (2009-03-20) and later if using hyperref.sty.
% The latest version can be obtained at:
% http://www.ctan.org/pkg/cite
% The documentation is contained in the cite.sty file itself.






% *** GRAPHICS RELATED PACKAGES ***
%
\ifCLASSINFOpdf
  % \usepackage[pdftex]{graphicx}
  % declare the path(s) where your graphic files are
  % \graphicspath{{../pdf/}{../jpeg/}}
  % and their extensions so you won't have to specify these with
  % every instance of \includegraphics
  % \DeclareGraphicsExtensions{.pdf,.jpeg,.png}
\else
  % or other class option (dvipsone, dvipdf, if not using dvips). graphicx
  % will default to the driver specified in the system graphics.cfg if no
  % driver is specified.
  % \usepackage[dvips]{graphicx}
  % declare the path(s) where your graphic files are
  % \graphicspath{{../eps/}}
  % and their extensions so you won't have to specify these with
  % every instance of \includegraphics
  % \DeclareGraphicsExtensions{.eps}
\fi
% graphicx was written by David Carlisle and Sebastian Rahtz. It is
% required if you want graphics, photos, etc. graphicx.sty is already
% installed on most LaTeX systems. The latest version and documentation
% can be obtained at: 
% http://www.ctan.org/pkg/graphicx
% Another good source of documentation is "Using Imported Graphics in
% LaTeX2e" by Keith Reckdahl which can be found at:
% http://www.ctan.org/pkg/epslatex
%
% latex, and pdflatex in dvi mode, support graphics in encapsulated
% postscript (.eps) format. pdflatex in pdf mode supports graphics
% in .pdf, .jpeg, .png and .mps (metapost) formats. Users should ensure
% that all non-photo figures use a vector format (.eps, .pdf, .mps) and
% not a bitmapped formats (.jpeg, .png). The IEEE frowns on bitmapped formats
% which can result in "jaggedy"/blurry rendering of lines and letters as
% well as large increases in file sizes.
%
% You can find documentation about the pdfTeX application at:
% http://www.tug.org/applications/pdftex





% *** MATH PACKAGES ***
%
%\usepackage{amsmath}
% A popular package from the American Mathematical Society that provides
% many useful and powerful commands for dealing with mathematics.
%
% Note that the amsmath package sets \interdisplaylinepenalty to 10000
% thus preventing page breaks from occurring within multiline equations. Use:
%\interdisplaylinepenalty=2500
% after loading amsmath to restore such page breaks as IEEEtran.cls normally
% does. amsmath.sty is already installed on most LaTeX systems. The latest
% version and documentation can be obtained at:
% http://www.ctan.org/pkg/amsmath





% *** SPECIALIZED LIST PACKAGES ***
%
%\usepackage{algorithmic}
% algorithmic.sty was written by Peter Williams and Rogerio Brito.
% This package provides an algorithmic environment fo describing algorithms.
% You can use the algorithmic environment in-text or within a figure
% environment to provide for a floating algorithm. Do NOT use the algorithm
% floating environment provided by algorithm.sty (by the same authors) or
% algorithm2e.sty (by Christophe Fiorio) as the IEEE does not use dedicated
% algorithm float types and packages that provide these will not provide
% correct IEEE style captions. The latest version and documentation of
% algorithmic.sty can be obtained at:
% http://www.ctan.org/pkg/algorithms
% Also of interest may be the (relatively newer and more customizable)
% algorithmicx.sty package by Szasz Janos:
% http://www.ctan.org/pkg/algorithmicx




% *** ALIGNMENT PACKAGES ***
%
%\usepackage{array}
% Frank Mittelbach's and David Carlisle's array.sty patches and improves
% the standard LaTeX2e array and tabular environments to provide better
% appearance and additional user controls. As the default LaTeX2e table
% generation code is lacking to the point of almost being broken with
% respect to the quality of the end results, all users are strongly
% advised to use an enhanced (at the very least that provided by array.sty)
% set of table tools. array.sty is already installed on most systems. The
% latest version and documentation can be obtained at:
% http://www.ctan.org/pkg/array


% IEEEtran contains the IEEEeqnarray family of commands that can be used to
% generate multiline equations as well as matrices, tables, etc., of high
% quality.




% *** SUBFIGURE PACKAGES ***
%\ifCLASSOPTIONcompsoc
%  \usepackage[caption=false,font=normalsize,labelfont=sf,textfont=sf]{subfig}
%\else
%  \usepackage[caption=false,font=footnotesize]{subfig}
%\fi
% subfig.sty, written by Steven Douglas Cochran, is the modern replacement
% for subfigure.sty, the latter of which is no longer maintained and is
% incompatible with some LaTeX packages including fixltx2e. However,
% subfig.sty requires and automatically loads Axel Sommerfeldt's caption.sty
% which will override IEEEtran.cls' handling of captions and this will result
% in non-IEEE style figure/table captions. To prevent this problem, be sure
% and invoke subfig.sty's "caption=false" package option (available since
% subfig.sty version 1.3, 2005/06/28) as this is will preserve IEEEtran.cls
% handling of captions.
% Note that the Computer Society format requires a larger sans serif font
% than the serif footnote size font used in traditional IEEE formatting
% and thus the need to invoke different subfig.sty package options depending
% on whether compsoc mode has been enabled.
%
% The latest version and documentation of subfig.sty can be obtained at:
% http://www.ctan.org/pkg/subfig




% *** FLOAT PACKAGES ***
%
%\usepackage{fixltx2e}
% fixltx2e, the successor to the earlier fix2col.sty, was written by
% Frank Mittelbach and David Carlisle. This package corrects a few problems
% in the LaTeX2e kernel, the most notable of which is that in current
% LaTeX2e releases, the ordering of single and double column floats is not
% guaranteed to be preserved. Thus, an unpatched LaTeX2e can allow a
% single column figure to be placed prior to an earlier double column
% figure.
% Be aware that LaTeX2e kernels dated 2015 and later have fixltx2e.sty's
% corrections already built into the system in which case a warning will
% be issued if an attempt is made to load fixltx2e.sty as it is no longer
% needed.
% The latest version and documentation can be found at:
% http://www.ctan.org/pkg/fixltx2e


%\usepackage{stfloats}
% stfloats.sty was written by Sigitas Tolusis. This package gives LaTeX2e
% the ability to do double column floats at the bottom of the page as well
% as the top. (e.g., "\begin{figure*}[!b]" is not normally possible in
% LaTeX2e). It also provides a command:
%\fnbelowfloat
% to enable the placement of footnotes below bottom floats (the standard
% LaTeX2e kernel puts them above bottom floats). This is an invasive package
% which rewrites many portions of the LaTeX2e float routines. It may not work
% with other packages that modify the LaTeX2e float routines. The latest
% version and documentation can be obtained at:
% http://www.ctan.org/pkg/stfloats
% Do not use the stfloats baselinefloat ability as the IEEE does not allow
% \baselineskip to stretch. Authors submitting work to the IEEE should note
% that the IEEE rarely uses double column equations and that authors should try
% to avoid such use. Do not be tempted to use the cuted.sty or midfloat.sty
% packages (also by Sigitas Tolusis) as the IEEE does not format its papers in
% such ways.
% Do not attempt to use stfloats with fixltx2e as they are incompatible.
% Instead, use Morten Hogholm'a dblfloatfix which combines the features
% of both fixltx2e and stfloats:
%
% \usepackage{dblfloatfix}
% The latest version can be found at:
% http://www.ctan.org/pkg/dblfloatfix




%\ifCLASSOPTIONcaptionsoff
%  \usepackage[nomarkers]{endfloat}
% \let\MYoriglatexcaption\caption
% \renewcommand{\caption}[2][\relax]{\MYoriglatexcaption[#2]{#2}}
%\fi
% endfloat.sty was written by James Darrell McCauley, Jeff Goldberg and 
% Axel Sommerfeldt. This package may be useful when used in conjunction with 
% IEEEtran.cls'  captionsoff option. Some IEEE journals/societies require that
% submissions have lists of figures/tables at the end of the paper and that
% figures/tables without any captions are placed on a page by themselves at
% the end of the document. If needed, the draftcls IEEEtran class option or
% \CLASSINPUTbaselinestretch interface can be used to increase the line
% spacing as well. Be sure and use the nomarkers option of endfloat to
% prevent endfloat from "marking" where the figures would have been placed
% in the text. The two hack lines of code above are a slight modification of
% that suggested by in the endfloat docs (section 8.4.1) to ensure that
% the full captions always appear in the list of figures/tables - even if
% the user used the short optional argument of \caption[]{}.
% IEEE papers do not typically make use of \caption[]'s optional argument,
% so this should not be an issue. A similar trick can be used to disable
% captions of packages such as subfig.sty that lack options to turn off
% the subcaptions:
% For subfig.sty:
% \let\MYorigsubfloat\subfloat
% \renewcommand{\subfloat}[2][\relax]{\MYorigsubfloat[]{#2}}
% However, the above trick will not work if both optional arguments of
% the \subfloat command are used. Furthermore, there needs to be a
% description of each subfigure *somewhere* and endfloat does not add
% subfigure captions to its list of figures. Thus, the best approach is to
% avoid the use of subfigure captions (many IEEE journals avoid them anyway)
% and instead reference/explain all the subfigures within the main caption.
% The latest version of endfloat.sty and its documentation can obtained at:
% http://www.ctan.org/pkg/endfloat
%
% The IEEEtran \ifCLASSOPTIONcaptionsoff conditional can also be used
% later in the document, say, to conditionally put the References on a 
% page by themselves.




% *** PDF, URL AND HYPERLINK PACKAGES ***
%
%\usepackage{url}
% url.sty was written by Donald Arseneau. It provides better support for
% handling and breaking URLs. url.sty is already installed on most LaTeX
% systems. The latest version and documentation can be obtained at:
% http://www.ctan.org/pkg/url
% Basically, \url{my_url_here}.




% *** Do not adjust lengths that control margins, column widths, etc. ***
% *** Do not use packages that alter fonts (such as pslatex).         ***
% There should be no need to do such things with IEEEtran.cls V1.6 and later.
% (Unless specifically asked to do so by the journal or conference you plan
% to submit to, of course. )


% correct bad hyphenation here
\hyphenation{op-tical net-works semi-conduc-tor}


\begin{document}
%
% paper title
% Titles are generally capitalized except for words such as a, an, and, as,
% at, but, by, for, in, nor, of, on, or, the, to and up, which are usually
% not capitalized unless they are the first or last word of the title.
% Linebreaks \\ can be used within to get better formatting as desired.
% Do not put math or special symbols in the title.
\title{Route Optimization in Smart Villages: \\ a Graph Neural Network Approach}
%
%
% author names and IEEE memberships
% note positions of commas and nonbreaking spaces ( ~ ) LaTeX will not break
% a structure at a ~ so this keeps an author's name from being broken across
% two lines.
% use \thanks{} to gain access to the first footnote area
% a separate \thanks must be used for each paragraph as LaTeX2e's \thanks
% was not built to handle multiple paragraphs
%

\author{Alberto Durán-López, Daniel Bolaños-Martinez, Zaid Almahmoud, Chandresh Pravin, Suparna De,~\IEEEmembership{Member,~IEEE,} and Maria Bermudez-Edo% <-this % stops a space
\thanks{A. Durán-López, D. Bolaños-Martinez and M. Bermudez-Edo are with the CITIC-UGR and IATUR institutes, University of Granada, Spain. (e-mail: \{albduranlopez, danibolanos, mbe\}@ugr.es)}% 
\thanks{Z. Almahmoud is with the Artificial Intelligence and Media Lab (AIM Lab), Northwestern University in Qatar, Qatar. (e-mail: zaid.almahmoud@northwestern.edu)}%<-this % stops a space
\thanks{C. Pravin and S. De are with the School of Computer Science and Electronic Engineering, University of Surrey, UK. (email: \{c.pravin, s.de\}@surrey.ac.uk)}% <-this % stops a space
}

% note the % following the last \IEEEmembership and also \thanks - 
% these prevent an unwanted space from occurring between the last author name
% and the end of the author line. i.e., if you had this:
% 
% \author{....lastname \thanks{...} \thanks{...} }
%                     ^------------^------------^----Do not want these spaces!
%
% a space would be appended to the last name and could cause every name on that
% line to be shifted left slightly. This is one of those "LaTeX things". For
% instance, "\textbf{A} \textbf{B}" will typeset as "A B" not "AB". To get
% "AB" then you have to do: "\textbf{A}\textbf{B}"
% \thanks is no different in this regard, so shield the last } of each \thanks
% that ends a line with a % and do not let a space in before the next \thanks.
% Spaces after \IEEEmembership other than the last one are OK (and needed) as
% you are supposed to have spaces between the names. For what it is worth,
% this is a minor point as most people would not even notice if the said evil
% space somehow managed to creep in.



% The paper headers
\markboth{IEEE Internet of Things Journal}%
{Durán-López \MakeLowercase{\textit{et al.}}: Route Optimization in Smart Villages}
% The only time the second header will appear is for the odd numbered pages
% after the title page when using the twoside option.
% 
% *** Note that you probably will NOT want to include the author's ***
% *** name in the headers of peer review papers.                   ***
% You can use \ifCLASSOPTIONpeerreview for conditional compilation here if
% you desire.




% If you want to put a publisher's ID mark on the page you can do it like
% this:
%\IEEEpubid{0000--0000/00\$00.00~\copyright~2015 IEEE}
% Remember, if you use this you must call \IEEEpubidadjcol in the second
% column for its text to clear the IEEEpubid mark.



% use for special paper notices
%\IEEEspecialpapernotice{(Invited Paper)}




% make the title area
\maketitle

% As a general rule, do not put math, special symbols or citations
% in the abstract or keywords.
\begin{comment}
\begin{abstract}
\textbf{Tourist visitation} plays an important role in smart village management, influencing resource allocation and congestion control. In this work, we leverage vehicle detection with Internet of Things (IoT) devices to capture real-world route visitation in smart villages. Although these sensors provide data in multiple formats, representing the information as graphs uncovers intricate spatial and temporal relationships that conventional tabular or time-series models often miss. We propose a framework that employs Graph Neural Networks (GNNs).   We collect each vehicle’s detections from License Plate Recognition cameras installed at four nodes across three villages, covering 2 years and 10 months and 467,773 distinct vehicles. We build a graph for each vehicle that maps its route across all cameras. Using these graphs, our study explores various GNN architectures, including spectral, spatial, recurrent, and attention-based models, with a special focus on incorporating edge features. We also apply explainability methods to identify the most influential graph components driving model predictions. To optimize performance, we integrate Neural Architecture Search and Hyperparameter Optimization (NAS-HPO) using a multivariate objective that maximizes the F1-score while minimizing computational time. Analysis of the results reveals that models using edge attributes perform better; specifically, the E-GAT model. We apply the NAS-HPO approach to this E-GAT base model, and the improved version achieves the highest performance across all metrics, demonstrating a 5.71\% improvement in F1-score and a 20.87\% decrease in execution time.
\end{abstract}
\end{comment}


\begin{abstract}
In the Internet-of-Things (IoT) domain, modeling the intricate relationships within spatiotemporal data from sensors is a significant challenge. Conventional methods, which often rely on tabular or time-series formats, frequently fail to capture the underlying complex dynamics. In contrast, graph-based representations excel at this task, but they have seldom been explored in tourism and remain unexplored in small villages, where IoT networks involve only a handful of sensors or nodes. We propose an explainable and optimized Graph Neural Network (GNN) methodology to predict return visits in smart villages. We collected data from License Plate Recognition cameras across three villages, tracking 467,773 distinct vehicles over 2 years and 11 months. For each vehicle, we construct a unique graph that maps its specific route across all cameras. We conduct a comprehensive study of various GNN architectures, including spectral, spatial, recurrent, and attention-based models, with a special focus on the impact of incorporating edge features. We also apply explainability methods to identify the most influential graph components driving model predictions. To optimize performance, we integrate Neural Architecture Search and Hyperparameter Optimization (NAS-HPO) using a multivariate objective that maximizes the F1-score while minimizing computational time. Analysis of the results reveals that models using edge attributes perform better; specifically, the E-GAT model. We apply the NAS-HPO approach to this E-GAT base model, and the improved version achieves the highest performance across all metrics, demonstrating a 5.71\% improvement in F1-score and a 20.87\% decrease in execution time.
\end{abstract}


% Note that keywords are not normally used for peerreview papers.
\begin{IEEEkeywords}
IoT, GNN, Explainability, NAS-HPO
\end{IEEEkeywords}






% For peer review papers, you can put extra information on the cover
% page as needed:
% \ifCLASSOPTIONpeerreview
% \begin{center} \bfseries EDICS Category: 3-BBND \end{center}
% \fi
%
% For peerreview papers, this IEEEtran command inserts a page break and
% creates the second title. It will be ignored for other modes.
\IEEEpeerreviewmaketitle



\section{Introduction} \label{sec:introduccion}
% The very first letter is a 2 line initial drop letter followed
% by the rest of the first word in caps.
% 
% form to use if the first word consists of a single letter:
% \IEEEPARstart{A}{demo} file is ....
% 
% form to use if you need the single drop letter followed by
% normal text (unknown if ever used by the IEEE):
% \IEEEPARstart{A}{}demo file is ....
% 
% Some journals put the first two words in caps:
% \IEEEPARstart{T}{his demo} file is ....
% 
% Here we have the typical use of a "T" for an initial drop letter
% and "HIS" in caps to complete the first word.
\IEEEPARstart{T}{he} expansion of Internet of Things (IoT) devices and accompanying sensor technologies has triggered an increase in the generation of daily data \cite{lo2022graphsage}. While traditionally this type of data has been organized in tabular or time-series formats \cite{Lakshmanna2022Review}, many real-world systems, ranging from sensor and transportation systems to social networks, are inherently interconnected and exhibit relationships that are most effectively represented as graphs \cite{li2024survey}. This graph-structured data poses new challenges for conventional machine learning (ML) approaches, which typically struggle to capture the detailed relational and structural information embedded within graphs. Graph Neural Networks (GNNs) \cite{kipf2016semi} have emerged as a powerful approach specifically developed to address these challenges \cite{wu2020comprehensive}. These models operate directly on graph data by aggregating and propagating features along the connections between nodes \cite{Zhou2020Review}. This process not only captures local interactions but also integrates global structural information, ensuring that the intricate relationships inherent in the data are effectively preserved and exploited \cite{Zheng2023TemporalGNN}.


When data is represented as a graph, it must be transformed into embeddings before being input into GNNs, a preprocessing step that significantly impacts downstream performance. Generating representative embeddings is necessary for robust outcomes in tasks such as graph classification, clustering, and node or edge prediction. These embeddings must not only preserve the structural and semantic information of the data but also enhance it by introducing additional representational dimensions that make the data more useful for models \cite{Xiao2023Heterogeneous}. Traditional GNNs, such as Graph Convolutional Networks (GCN) \cite{kipf2016semi}, Gated Graph Convolutional Networks (GGCN) \cite{li2015gated}, and Graph Attention Networks (GAT) \cite{velivckovic2017graph}, have primarily focused on capturing node features and overall graph structure while often overlooking the information encoded within edge features \cite{gong2019exploiting}. Recent advances have addressed this limitation by modifying the message passing process to explicitly incorporate edge attributes, as seen in models such as the Relational Graph Convolutional Network (R-GCN) \cite{schlichtkrull2018modeling} and various heterogeneous GNN architectures \cite{Zhu2019Relation}. This integration of edge information not only enhances model performance but also works to improve explainability, offering valuable insights into how specific nodes and edges contribute to the decisions made by models \cite{Yuan2023Explainability}.


Over recent years, new GNNs have been developed to address the intricate and domain-specific challenges posed by different types of graph data \cite{li2024survey}. These architectures include spectral convolutional \cite{defferrard2016convolutional}, spatial \cite{hamilton2017inductive}, recurrent \cite{li2015gated}, and attention-based models \cite{velivckovic2017graph}. Additionally, spatio-temporal models \cite{huang2024dst} have been introduced to capture changes in graph structures over time. However, performance variations between these models are often driven by the inherent characteristics of the data more than by the models themselves \cite{Zhou2020Review}.
Moreover, certain GNN architectures function as \textit{black boxes}, making it difficult to understand how predictions are made. Consequently, identifying the nodes and edges that most significantly influence predictions is important for understanding model behavior and ensuring robustness in real-world applications \cite{Zhang2024TrustworthyGNNs}. Explanation techniques, such as GNNExplainer \cite{ying2019gnnexplainer}, shed light on the key features and relationships driving model performance, thus improving transparency and facilitating validation.


Despite the advancements in new GNN architectures, a fundamental challenge persists: finding an effective model configuration for a specific graph dataset. Neural Architecture Search (NAS) \cite{gao2019graphnasgraphneuralarchitecture} has emerged as a solution to this challenge by systematically exploring various configurations of layers and hyperparameters. This automated approach is in response to the observation that performance variations of GNNs can often be attributed more to the characteristics of the data than model selection \cite{Zhou2020Review}. By leveraging NAS, researchers can efficiently identify optimal architectures tailored to specific graph tasks \cite{zhang2019heterogeneous}. This not only accelerates development but also enhances performance by discovering architectures suited to the intrinsic properties of each dataset \cite{zoph2017neural}.


In real-world traffic and transportation applications, road networks are commonly represented as graphs, where roads correspond to edges and intersections to nodes, thereby facilitating further analysis \cite{jana2023edge}. This graph-based approach has proven to be effective for various tasks, such as route planning and traffic prediction \cite{li2024survey}. In the broader smart tourism domain, machine learning has been used to optimize tourist itineraries by modeling attractions and paths as a graph network \cite{mangini2021novel}, as well as to predict user behavior on online travel platforms for hotel recommendation \cite{li2023multitask}.

%
% However, the smart village domain remains under-explored in this regard. 
%
However, the smart village domain remains under-explored, particularly in applications of GNNs for mobility analysis, where the resulting graphs tend to be small, raising questions about the effectiveness of approaches on transportation graphs with a small number of nodes.
%

To address this gap, we present a methodology that leverages graph-based ML with explainability on a real smart village traffic dataset. The dataset, collected via IoT cameras with automatic number plate recognition (ANPR), is used to construct a road network graph reflecting vehicle movements. Using this data, we predict repeat tourist visits, yielding societal benefits by reducing overcrowding \cite{bichler2021designing}, economic benefits through steady revenue \cite{acharya2023linking}, and planning benefits by guiding infrastructure and resource allocation \cite{alfarhan2022first}. By applying GNN models along with explainability techniques, we can understand which features most influence traffic patterns in the village, turning black-box predictions into actionable insights. While NAS has been employed in other domains to automatically improve deep neural networks (including GNNs \cite{gao2021graph}), to our knowledge, it has never been applied to tourist mobility classification in a smart village context. This work is therefore novel in combining graph-based route modeling, IoT data, GNN explainability, and an exploration of NAS to optimize models for smart village traffic management, addressing an important gap in the current research.

This study addresses these gaps by proposing an integrated framework that combines IoT-based graph construction, GNN explainability, and automated GNN architecture search for rural mobility analysis. Our specific objectives are: first, to evaluate the predictive performance of GNN models on small-scale transportation networks; second, to assess the explainability of GNN predictions for tourism; and third, to determine the efficiency improvements of using automated neural architecture search.


The main contributions of our approach include:
\begin{itemize}
    \item We present a methodology and validate it on our case study dataset to forecast tourism in a rural smart village. Our approach transforms real-world smart village data into graph representations and employs explainability techniques on the GNN. By evaluating both node and edge feature importance, we gain clear insights into the decision-making process.
    \item We implement a multivariate NAS and Hyperparameter Optimization (NAS-HPO) strategy to explore various layer configurations and hyperparameter settings, aiming to maximize the F1-score while reducing execution time.
    \item We construct a real-world dataset from a road network using IoT ANPR sensor data, covering 2 years and 11 months of vehicle routes in a smart village. This dataset, comprising 467,773 distinct vehicles and preserving the sequential patterns of tourist routes, enables thorough evaluation of our approach.
\end{itemize}


%IEEE IoT journal uses remainder 
The remaining sections of this article are organized as follows. Section \ref{sec:related_work} provides a review of the existing literature on GNN models and state-of-the-art on smart village case studies. The proposed methodology is outlined in Section \ref{sec:methodology}. Section \ref{sec:experiments} details the experiments conducted on our smart village case study dataset, with results analyzed in Section \ref{sec:results_discussion}. Finally, Section \ref{sec:conclusion} concludes the paper.

\section{Related work} \label{sec:related_work}


In this section, we provide a comprehensive review of GNN literature, covering spectral, spatial, attention-based networks, and other related architectures. In addition to foundational and theoretical works, we also incorporate recent empirical studies on tourism and smart villages relevant to our work, while highlighting our contribution.

\subsection{Spectral Graph Neural Networks}


Spectral GNNs leverage graph signal processing techniques to perform convolutions in the spectral domain, effectively capturing structural information from graph-structured data. ChebNet \cite{defferrard2016convolutional} introduced an efficient formulation of spectral convolutions using Chebyshev polynomial approximations, enabling fast and localized graph filtering with linear computational complexity. Building on this foundation, Kipf and Welling \cite{kipf2016semi} proposed GCN, a first-order approximation of spectral convolutions that enables scalable semi-supervised learning by efficiently aggregating node features using a normalized Laplacian. However, GCN does not fully exploit edge attributes, which can encode essential relational information. To address this, E-GCN \cite{gong2019exploiting} extends the GCN framework by integrating edge features through a doubly stochastic normalization strategy and multi-dimensional edge-aware transformations, leading to enhanced representation learning for tasks such as node and graph classification.

\subsection{Spatial Graph Neural Networks}

Spatial GNNs operate directly on a graph’s structure by defining convolutional operations in the spatial domain, typically aggregating information from local node neighborhoods. One of the most well-known spatial GNNs, GraphSAGE \cite{hamilton2017inductive}, provides an inductive learning framework that generates node embeddings by sampling and aggregating features from neighboring nodes. This approach enables generalization to unseen nodes and scales efficiently to large graphs. Building on GraphSAGE, E-GraphSAGE \cite{lo2022graphsage} extends GraphSAGE by incorporating edge features to modulate the message passing based on relationship information, demonstrating superior performance on benchmark datasets.

Meanwhile, Message Passing Neural Networks (MPNNs) \cite{gilmer2017neural} unify various GNN models under a common framework and have been particularly effective in molecular property prediction tasks. In the context of point clouds, Dynamic Graph CNN (DGCNN) \cite{wang2019dynamic} introduces EdgeConv, a module that dynamically computes graph structures at each layer to capture both local and global shape properties, significantly improving classification and segmentation tasks.  

For knowledge graphs, relational GCNs (R-GCNs) \cite{schlichtkrull2018modeling} extend GCNs to handle multi-relational data by incorporating relation-specific transformations, proving effective in link prediction and entity classification. Furthermore, heterogeneous graph neural networks (HetGNNs) \cite{zhang2019heterogeneous} address the challenges of learning representations in heterogeneous graphs by leveraging a random walk sampling strategy and a two-stage aggregation mechanism, achieving strong performance in tasks like recommendation and node classification.



\subsection{Attention-based Graph Neural Networks}

Attention-based GNNs enhance message passing by dynamically weighting node interactions, allowing models to differentiate the importance of neighboring nodes. 

Graph Attention Networks (GATs) \cite{velivckovic2017graph} introduce masked self-attentional layers to assign different importance scores to neighboring nodes, improving upon traditional graph convolutional approaches. By leveraging multi-head attention, GAT enables efficient learning on both transductive and inductive tasks while avoiding expensive matrix operations. The model achieves state-of-the-art results on various benchmarks.

Expanding on this, Edge-featured Graph Attention Networks (E-GATs) \cite{wang2021egat, chen2021edge} incorporate edge attributes into the attention mechanism, addressing limitations in node-centric GNNs. E-GAT learns representations by integrating edge features into attention weight calculations and updating edge attributes alongside node embeddings. This approach enhances performance on edge-sensitive graph tasks, making E-GAT particularly effective for applications in domains like financial transaction networks.



\subsection{Graph Autoencoders and Recurrent Graph Neural Networks}

Graph Autoencoders (GAEs) are unsupervised learning models designed to encode graph-structured data into latent representations and reconstruct the graph topology from them. One of the most notable variants, Variational Graph Autoencoder (VGAE) \cite{kipf2016variational}, extends the variational autoencoder (VAE) framework to graphs by leveraging a graph convolutional network (GCN) as the encoder and an inner-product decoder. By incorporating node features into the latent variable model, VGAE achieves competitive performance in unsupervised link prediction tasks on benchmark citation network datasets.

Recurrent GNNs extend GNN architectures by incorporating recurrent mechanisms, allowing iterative updates of node embeddings through multiple message-passing steps. One such model, Gated Graph Neural Networks (GGNN) \cite{li2015gated}, improves upon graph neural networks by integrating gated recurrent units (GRUs) to iteratively refine node representations based on their local neighborhoods. GGNNs have demonstrated strong performance in various structured learning tasks, including program verification and algorithm learning.

Building upon GGNN, Edge-Gated Graph Neural Networks (E-GGNN) \cite{zheng2024edge} enhance message passing by introducing edge-aware gating mechanisms. These mechanisms dynamically regulate information flow along edges, allowing the model to capture more nuanced relational structures. E-GGNN has shown effectiveness in specialized domains, such as few-shot learning and molecular property prediction, where modeling edge information is critical for improving classification and regression tasks.


\subsection{Spatio-Temporal Graph Neural Networks}

Spatio-Temporal GNNs extend conventional GNNs by modeling temporal dependencies in addition to spatial dependencies, making them effective for dynamic graph-based tasks such as traffic forecasting and human activity recognition. These models simultaneously leverage spatial and temporal information to capture evolving patterns in structured data.

Spatio-Temporal Graph Convolutional Network (STGCN) \cite{yu2017spatio} integrates spatial graph convolutions with temporal convolutional networks to effectively model dynamic traffic patterns over time. Instead of relying on standard recurrent or convolutional layers, STGCN processes traffic data as a graph structure, enabling more efficient and accurate forecasting. By modeling multi-scale traffic networks, STGCN significantly improves predictive performance on real-world datasets.

Building upon spatio-temporal modeling, Dynamic Spatio-Temporal Graph Transformer Network (DST-GTN) \cite{huang2024dst} introduces a transformer-based approach that dynamically adapts to changing spatial relationships over time. Unlike traditional methods that treat spatial structures as static, DST-GTN captures dynamic adjacency relationships and refines both global and local spatio-temporal representations using adaptive filters. This enables the model to extract rich spatio-temporal features from time-series traffic data, achieving state-of-the-art results in various forecasting tasks.

Recently, a Bayesian approximation \cite{almahmoud2025forecasting} of a Multivariate Time-series Graph Neural Network (B-MTGNN) model \cite{wu2020connecting} was proposed for forecasting multivariate time-series data represented as a graph. This approach models spatial and temporal dependencies and adaptively learns hidden relationships between different variables to enhance multivariate time-series prediction, particularly for complex datasets such as cybersecurity trends. Additionally, the model quantifies epistemic uncertainty using the Monte Carlo dropout method \cite{almahmoud2023holistic}, improving reliability through uncertainty estimation.



\subsection{Hierarchical Pooling Graph Neural Networks}

Hierarchical Pooling GNNs address the limitations of traditional GNN architectures by enabling the learning of hierarchical representations of graphs. This capability is particularly beneficial for graph classification tasks, where a comprehensive understanding of the graph's structure is important.

DiffPool \cite{ying2018hierarchical} introduces a differentiable graph pooling module that generates hierarchical representations. By learning soft cluster assignments for nodes at each layer, DiffPool coarsens graphs and creates fixed-size representations suitable for classification. This method can be integrated with various GNN architectures in an end-to-end manner, improving performance on graph classification benchmarks by 5-10\% compared to existing pooling methods, and achieving state-of-the-art results on multiple datasets.


Building on the concept of hierarchical representation learning, Graph U-Nets \cite{gao2019graph} introduce novel graph pooling (gPool) and unpooling (gUnpool) operations to address the unique challenges posed by graph data. The gPool layer selectively chooses nodes based on their scalar projection values onto a trainable vector, effectively coarsening the graph for downstream classification tasks by selecting the highest scoring nodes (TopKPool). The inverse gUnpool layer reconstructs the original graph structure from the pooled representation. The introduction of these operations enables the development of an encoder-decoder architecture tailored for graph data, demonstrating superior performance in node and graph classification tasks. Table \ref{tab:gnn_overview_final} summarizes the related work on GNNs. 


\begin{table*}[htbp]
  \centering
  \caption{Overview of Graph Neural Networks}
  \label{tab:gnn_overview_final}
  \renewcommand{\arraystretch}{1.2}
  \begin{adjustbox}{max width=\textwidth}
    \begin{tabular}{p{2cm}lp{9cm}p{3.5cm}}
      \toprule
      \textbf{Category} & \textbf{Name} & \textbf{Methodology} & \textbf{Application/scope} 
      \\
      \midrule
      \multirow{3}{*}{Spectral} 
        & ChebNet \cite{defferrard2016convolutional}   & Approximates spectral filters using Chebyshev polynomials for localized graph convolution. & Node classification \\ 
        & GCN \cite{kipf2016semi} & Uses a first-order approximation of spectral convolutions for scalable feature aggregation. & Node classification  \\ 
        & E-GCN \cite{gong2019exploiting}& Extends GCN by integrating edge features with a multi-dimensional edge-aware strategy. & Node and graph classification and regression  \\ 
      \midrule
      \multirow{6}{*}{Spatial} 
        & GraphSAGE \cite{hamilton2017inductive} & Learns node embeddings by sampling and aggregating neighbor features. & Node classification \\ 
        & E-GraphSAGE \cite{lo2022graphsage} & Incorporates edge features into GraphSAGE for edge-aware aggregation. & Node and graph classification  \\ 
        & MPNN \cite{gilmer2017neural} & General message-passing framework that updates node embeddings for diverse applications. & Molecular property prediction  \\ 
        & EdgeConv \cite{wang2019dynamic} & Dynamically computes edge features for geometric data learning. & Point cloud classification and segmentation \\ 
        & R-GCN \cite{schlichtkrull2018modeling} & Extends GCN to multi-relational graphs with relation-specific transformations. & Link prediction, entity classification \\ 
        & HetGNN \cite{zhang2019heterogeneous}  & Aggregates heterogeneous graph data using type-specific strategies. & Link prediction, node classification and clustering \\ 
      \midrule
      \multirow{2}{*}{Attention} 
        & GAT \cite{velivckovic2017graph} & Uses self-attention to weight neighbor contributions in message passing. & Node classification  \\ 
        & E-GAT \cite{wang2021egat, chen2021edge} & Integrates edge features into attention mechanisms for better representation. & Node and graph classification \\ 
      \midrule
      \multirow{1}{*}{Autoencoders} 
        & VGAE \cite{kipf2016variational} & Variational autoencoder for unsupervised graph representation learning. & Link prediction  \\ 
      \midrule
      \multirow{2}{*}{Recurrent} 
        & GGNN \cite{li2015gated} & Uses gated recurrent units (GRUs) to refine node representations iteratively. & Structured learning  \\ 
        & E-GGNN \cite{zheng2024edge,jiao2021edge} & Adds edge-aware gating to GGNN for better relational modeling. & Few-shot learning, molecular property prediction \\ 
      \midrule
      \multirow{3}{*}{Spatio-Temporal} 
        & STGCN \cite{yu2017spatio} & Combines spatial graph convolutions with temporal convolutions for time-series graphs. & Traffic forecasting  \\ 
        & DST-GTN \cite{huang2024dst} & Uses transformers for dynamic spatial-temporal relationship modeling. & Traffic forecasting \\ 
        & B-MTGNN \cite{almahmoud2025forecasting} & Bayesian model for multivariate time-series forecasting with uncertainty quantification. & Multivariate time-series forecasting  \\ 
      \midrule
      \multirow{2}{*}{\makecell{Hierarchical\\Pooling}} 
        & DiffPool \cite{ying2018hierarchical} & Learns hierarchical graph representations via differentiable pooling. & Graph classification \\ 
        & TopKPool \cite{gao2019graph} & Selects top-ranked nodes for hierarchical pooling and classification. & Node and graph classification \\ 
      \bottomrule
    \end{tabular}
  \end{adjustbox}
\end{table*}






\subsection{Comparative Analysis of Recent Works}

Recent studies have increasingly applied machine learning and deep learning techniques to tourism forecasting and smart village analytics. These efforts span a variety of approaches, from classical time-series models to graph-based and attention-based neural architectures. In this section, we review recent empirical work in two main areas relevant to our contribution: tourism forecasting and smart village systems. We end this section by positioning our contribution within this evolving landscape.


\subsubsection{Tourism Forecasting}


Some researchers employed classical time-series models to forecast tourism demand \cite{nunez2024machine}. Others integrated online behaviour and internet-search trends into prediction models with kernel-based and mixed-frequency learning approaches \cite{hu2025tourism}. The study in \cite{li2024forecasting} applied RobustSTL to decompose the series into a trend component and a residual component capturing both seasonal cycles and irregular fluctuations. It forecasted the trend with ARIMA, modeled the residual with an LSTM, and merged the two forecasts to produce the final prediction. \cite{chen2024forecasting} picked Baidu keywords via Boruta, added visitor counts, and used a CNN-BiLSTM framework to forecast monthly flows, outperforming both traditional and single-network models across three Chinese attractions.
The study in \cite{kim2021daily} focuses on daily inbound tourist volume in South Korea using a Multi-Head Attention CNN framework that integrates temporal attention to deliver more accurate predictions than other deep-learning approaches. Transformer-based models have been used in \cite{li2024enhancing}, where temporal fusion transformers are combined with Bayesian optimization to produce robust forecasts. Extensions of the transformer architecture appear in \cite{chen2025tourism}, addressing spatial–temporal dependencies and interpretability through modular fusion networks and masked attention mechanisms.


Recent work adopted graph-based architectures: the study in \cite{liang2025tourism} proposed a GNN exploiting spatiotemporal correlations with dual correlation matrices and attention mechanisms, achieving high accuracy across Chinese tourist cities. A graph-guided network modeling lag effects by dynamically encoding historical data into a bipartite graph appears in \cite{fan2023graph}. \cite{li2022tourism} uses a hybrid GCN–LSTM model to forecast monthly tourist arrivals across European countries and hourly visits to Beijing attractions, demonstrating that incorporating spatial effects through graph structures significantly reduces forecasting errors. The approach in \cite{wang2021multi} employs a multi-graph convolutional network to capture diverse relationships among attractions via multiple graph representations, outperforming baseline methods in tourist flow forecasting. \cite{hajisafi2023learning} introduces Busyness GNN (BysGNN) to predict POI visits in the United States by constructing dynamic graphs from temporal, geographic, and contextual features, yielding significant performance gains over state-of-the-art models.



\subsubsection{Smart Villages}

Data from satisfaction questionnaires have been employed to train and optimize genetic algorithms to develop an accommodation-recommendation system for smart villages \cite{suanpang2024can}. Similarly, \cite{surya2024tourism} proposes an integrated tourism model for smart villages and validates it via Partial Least Squares Structural Equation Modeling (PLS-SEM). Only a few works exploit sensor data within smart villages. For instance, one study leverages ANPR data, with socio-economic and seasonal features, to predict overnight stays in smart villages \cite{bolanos2024predicting}. Similarly, \cite{bolanos2025sasd} combines questionnaires and contextual datasets with ANPR feeds to forecast tourists’ intent to visit over the next 12 months.


\subsubsection{Research Contribution}

While tourism studies offer valuable contributions, they primarily focus on forecasting aggregated tourist volume or demand across cities or attractions, often using either time-series or structured input representations. Among the reviewed graph approaches, none evaluate multiple GNN models to determine which architecture works best for the task. They also do not use explainability to reveal how edge attributes contribute to predictions, nor do they apply NAS HPO tuning to jointly optimise model structure and training settings. Moreover, existing methods depend on aggregate flow statistics or static graphs built from cumulative data rather than on the individual time ordered paths of each traveller. This limitation is especially problematic in smart village settings where only a small number of sensor nodes are available and every single route matters.

%{\color{blue}While recent tourism research has embraced various machine learning techniques, it still focuses on predicting overall visitor counts, whether at the city, regional or attraction level, using time series data, search engine trends or static feature sets. These approaches rarely capture the finer details of individual travel routes or leverage real time sensor feeds to model repeat visits at the micro level.}

Existing smart village studies use contextual databases with limited accessibility (such as questionnaires), which restricts model training. In contrast, our work focuses on repeat visitation prediction in smart village environments by modeling real-world route data as graphs derived from IoT-based vehicle detection, without relying on additional information. We explore a diverse set of GNN architectures, including spectral, spatial, recurrent, and attention-based models, while also incorporating edge features. Furthermore, we employ NAS-HPO to jointly optimize performance and efficiency, addressing scalability and deployment in rural sensor-based systems. We additionally apply explainability methods to uncover graph structures driving predictions, highlighting a novel intersection between predictive modeling, efficient inference, and model interpretability in micro-level tourism prediction.


\section{Proposed Methodology} \label{sec:methodology}

Our methodology (see Figure \ref{fig:pipeline}) begins with the creation of a graph object for each dataset instance. 
This involves identifying which nodes and edges (with their attributes) form the graphs and converting them into graph objects (Section \ref{sec:graph_representation}). Next, we apply different GNN models to evaluate their performance in classifying these graphs. This step involves using a wide range of GNNs to determine which models adapt best to the data, and whether edge features provide additional information or performance improvements (Section \ref{sec:gnn_subsec}). We do this because GNNs often perform differently across datasets, showing that comparisons are task-specific \cite{errica2019fair}. We then use explainability techniques to understand feature importance in our best-performing model (Section \ref{sec:explainability}). This allows us to understand which nodes or edges are most influential and to determine the key drivers of the model’s predictions. Finally, after identifying the GNN model that best fits the classification task and using it as a base, we employ NAS-HPO to identify the optimal combination of layers and hyperparameters for the model’s architecture (Section \ref{sec:nas_subsec}).

\begin{figure*}[htbp]
    \centering
    \includegraphics[width=1\textwidth]{Images/pipeline.png}
    \caption{Pipeline and architecture of our framework}
    \label{fig:pipeline}
\end{figure*}


\subsection{Graphs Preprocessing and Representation} \label{sec:graph_representation}

If the initial dataset is not given as a graph object, it must be preprocessed into a graph structure. This involves defining which features become the nodes, which relationships become the edges, and whether nodes and edges carry additional attributes.

%The initial dataset is transformed into graph objects. 
We define a graph as $G = (V, E, X_v, X_e)$ (see \cite{gilmer2017neural,bora2023edge}), where:
\begin{itemize}
    \item $V$ is the set of nodes.
    \item $E \subseteq V \times V$ is the set of edges.
    \item $X_v \in \mathbb{R}^{|V| \times d}$ is the node feature matrix, with each node $v \in V$ associated with a feature vector $x_v \in \mathbb{R}^d$.
    \item $X_e \in \mathbb{R}^{|E| \times p}$ is the edge feature matrix, where each edge $e_{ij} \in E$, connecting nodes $v_i$ and $v_j$, is associated with a feature vector $x_{e_{ij}} \in \mathbb{R}^p$.
\end{itemize}

This transforms the problem of processing tabular or sequential camera detection data into a graph network problem, which is needed to structure the data in a way that GNN models can directly process. In our work, the graph structure is fixed for each instance, meaning that the nodes, edges, and their corresponding features remain constant over time.

\subsection{Graph Neural Network} \label{sec:gnn_subsec}

GNNs aim to learn expressive representations of nodes by iteratively aggregating and transforming information from their local neighborhoods. In the traditional formulation, the focus is on leveraging the graph structure and node features alone \cite{kipf2016semi,velivckovic2017graph}. For each node $v \in V$, let $h_v^{(l)}$ denote its hidden representation at layer $l$, with the initial representation given by $h_v^{(0)} = x_v$. The update mechanism at layer $l$ is defined as:
\begin{equation}
\label{eq:mp_traditional}
    \begin{aligned}
    m_v^{(l)} &= \operatorname{AGG}^{(l)}\Bigl(\{\, \phi^{(l)}\bigl(h_v^{(l-1)}, h_u^{(l-1)}\bigr) : u \in \mathcal{N}(v)\,\}\Bigr),\\[1mm]
    h_v^{(l)} &= \psi^{(l)}\bigl(h_v^{(l-1)}, m_v^{(l)}\bigr),
    \end{aligned}
\end{equation}

where:
\begin{itemize}
    \item $\mathcal{N}(v)$ denotes the set of neighbors of node $v$.
    \item $\phi^{(l)}$ is a message function that computes the contribution from neighbor $u$ to node $v$.
    \item $\operatorname{AGG}^{(l)}$ is an aggregation function (e.g., sum, mean, or max) that combines the messages from all neighbors.
    \item $\psi^{(l)}$ is an update function that fuses the previous representation of $v$ with the aggregated message.
\end{itemize}

After $L$ layers, each node’s final representation $h_v^{(L)}$ encapsulates information from its $L$-hop neighborhood. For graph-level tasks, a readout function $R(\cdot)$ aggregates these node representations into a single graph embedding:
\begin{equation}
    h_G = R\bigl(\{h_v^{(L)} : v \in V\}\bigr).
\end{equation}

While the framework above leverages only node features and the underlying graph structure, modern GNNs extend the message-passing process to incorporate edge features as well. In these models, the message function is modified to include edge attributes $x_{e_{uv}}$ from the graph $G = (V, E, X_v, X_e)$, leading to the following update:

\begin{equation}
\label{eq:mp_extended}
    \begin{aligned}
    m_v^{(l)} &= \operatorname{AGG}^{(l)}\Bigl(\{\, \phi^{(l)}\bigl(h_v^{(l-1)}, h_u^{(l-1)}, x_{e_{uv}}\bigr) : u \in \mathcal{N}(v)\,\}\Bigr),\\[1mm]
    h_v^{(l)} &= \psi^{(l)}\bigl(h_v^{(l-1)}, m_v^{(l)}\bigr).
    \end{aligned}
\end{equation}

In Equation \ref{eq:mp_extended}, while the message function $\phi^{(l)}$ has been adapted as detailed above to include edge features $x_{e_{uv}}$, the remaining components, specifically the aggregation function $\operatorname{AGG}^{(l)}$, the update function $\psi^{(l)}$, and the node representations $h_v^{(l-1)}$ and $h_u^{(l-1)}$, preserve their definitions and operational roles from the traditional formulation (Equation \ref{eq:mp_traditional}).

In our study, we aim to evaluate whether the inclusion of edge features $x_{e_{uv}}$ enables the network to capture additional relational and contextual information between nodes. To this end, we first employ traditional GNN models that rely solely on node features and the graph structure. Subsequently, we experiment with GNN architectures that incorporate edge features, as in Equation \ref{eq:mp_extended}. This comparative analysis will help us determine if integrating edge attributes improves the models’ ability to learn more detailed representations of complex graph-structured data.



\subsection{Explainability} \label{sec:explainability}

To interpret the predictions of our GNN models, we use an explainability method that shows the importance of both node features and edges. We use GNNExplainer \cite{ying2019gnnexplainer} because it produces masks for nodes and edges, which helps us understand what parts of the graph drive the prediction. This approach gives us insight into how the network works and whether including edge features improves the learned representations.

Given a graph $G = (V, E, X_v, X_e)$ and a trained GNN model $f$, the explainer produces two importance masks, the mask for node features $M_v$ and the mask for edges $M_e$:
\[
M_v \in [0,1]^{|V|\times d} \quad \text{and} \quad M_e \in [0,1]^{|E|},
\]
which indicate how relevant each node and edge feature is to the model’s prediction. Specifically, GNNExplainer introduces soft masks $X_v$ and $X_e$ to modulate the original inputs:
\[
\tilde{X}_v = X_v \odot M_v \quad \text{and} \quad \tilde{X}_e = X_e \odot M_e,
\]
where $\odot$ denotes element-wise multiplication. The perturbed graph is then $\tilde{G} = (V, E, \tilde{X}_v, \tilde{X}_e)$. The explainer learns these masks by minimizing the following objective:
\[
\min_{M_v, M_e} \; \mathcal{L}_{\text{pred}}\bigl(f(G), f(\tilde{G})\bigr) \;+\; \lambda_1 \|M_v\|_1 \;+\; \lambda_2 \|M_e\|_1,
\]
subject to $M_v \in [0,1]^{|V|\times d}$ and $M_e \in [0,1]^{|E|}$. Here:
\begin{itemize}
    \item The prediction loss $\mathcal{L}_{\text{pred}}\bigl(f(G), f(\tilde{G})\bigr)$ measures the discrepancy between the original model output $f(G)$ and the output on the perturbed graph $f(\tilde{G})$. We use cross-entropy loss because our task is a classification problem. %mse is for regression
    \item The $L_1$ norms $\|M_v\|_1$ and $\|M_e\|_1$ encourage sparsity in the masks, making the explanations simpler and easier to interpret.
    \item The hyperparameters $\lambda_1$ and $\lambda_2$ balance the trade-off between fidelity and sparsity in the learned masks.
\end{itemize}

By optimizing this objective, GNNExplainer identifies a subgraph and subset of node features that most strongly influence the model's prediction, thus providing interpretable insights into how the GNN processes graph-structured data. This process not only offers explainability into the decision-making process of the GNN but also informs the choice of model architecture by quantifying the contribution of node and edge information to overall performance.




\subsection{Neural Architecture Search} \label{sec:nas_subsec}

We perform NAS-HPO to design an effective GNN architecture \cite{izquierdo2021bag}. Our approach searches over a candidate configuration space $\Theta$, which consists of various combinations of network layers and hyperparameter settings. 

For the architectural hyperparameters, we allow 1–3 node and edge projection layers, where a single layer applies a linear transformation and using 2–3 layers enables richer, non‐linear representations. We also consider 1–4 GNN layers (e.g., E-GAT), since each layer performs one message-passing step by attending to neighbor information (and including edge features when available). Stacking more than four layers can lead to oversmoothing, so it is good practice to add residual connections for stability \cite{lell2023split}. We set fusion and message MLP depths to 1–3 layers \cite{stark20223d}, with an option of 0 message MLP layers to disable the message transformation entirely. To combine neighbor messages into a single vector before updating a node’s representation, we use mean, sum, and max aggregators \cite{you2020design}. Finally, for activation functions, which allow GNNs and their internal MLPs to learn complex mappings beyond affine combinations, we include ReLU, LeakyReLU, and GELU as widely used options \cite{rasool2024optimizing}.


For each configuration $\theta \in \Theta$, we train a GNN model $f_\theta$ with parameters $W$ by minimizing the training loss:
\[
W^*(\theta)=\arg\min_{W}\mathcal{L}_{\text{train}}(W,\theta).
\]
After training, we evaluate each configuration on a validation set by measuring the F1-score $F1(\theta,W^*(\theta))$ and recording the training time $T(\theta)$. Our goal is to achieve both high predictive performance and low computational cost; hence, we formulate a multivariate optimization problem by simultaneously minimizing the negative F1-score (so that a higher F1-score corresponds to a lower objective value) and the training time. This is expressed by finding the Pareto front:
\[
\theta^*=\operatorname{ParetoFront}\Bigl\{ \bigl(-F1(\theta,W^*(\theta)),\,T(\theta)\bigr) : \theta\in\Theta\Bigr\}.
\]
To efficiently explore the configuration space, we employ the Tree-Structured Parzen Estimator (TPE) as implemented in the TPASampler \cite{joy2022comprehensive}. TPE partitions the space based on a threshold $y^*$ and estimates two conditional densities:
\[
l(\theta)=p\bigl(\theta\mid f(\theta)<y^*\bigr) \quad \text{and} \quad g(\theta)=p\bigl(\theta\mid f(\theta)\ge y^*\bigr).
\]
%where $f(\theta)$ represents the combined performance metric. 
Here, $l(\theta)$ is the probability density function $\theta$ given that its performance metric $f(\theta)$ is good; and $g(\theta)$ is the probability density function of configurations $\theta$ given that their performance $f(\theta)$ is not good. The next candidate configuration $\theta_{\text{new}}$ is then selected by maximizing the ratio:
\[
\theta_{\text{new}}=\arg\max_{\theta\in\Theta}\frac{l(\theta)}{g(\theta)},
\]
which serves as a surrogate for the expected improvement, directing the search toward configurations likely to yield higher F1-score with lower training times.




\subsection{Model Training and Evaluation} \label{sec:model_training_eval}

Each GNN has its own characteristics and advantages. By adapting and testing different GNNs, we can evaluate which algorithms perform best for our problem, helping us identify the most suitable approach for the data. 

For the training hyperparameters, we consider the following ranges and options: the hidden dimension, representing each node’s intermediate embedding size, varies between 32 and 256, which covers the spot for many graph datasets \cite{song2021network}. We apply dropout with rates up to 0.5 during training to randomly mask node activations and include batch normalization at each layer to stabilize hidden activations and reduce internal covariate shift, accelerating convergence, since this combination markedly improves graph classification performance \cite{luobeyond}. We set the learning rate to range from $10^{-4}$ to $10^{-2}$, a common default for GNN training \cite{gu2025neubm}. Finally, we include undersampling as an on/off option to balance classes \cite{ma2025class}.

It is important to select various metrics to compare algorithm performance on test data, such as F1-score, recall, and precision \cite{naidu2023review}. Recall measures the proportion of actual positive cases that are correctly identified, highlighting the model’s ability to detect all relevant instances. Precision  quantifies the proportion of predicted positive cases that are truly positive, which is important for evaluating how many false positives occur. The F1-score, as the harmonic mean of recall and precision, provides a balanced measure that accounts for both false positives and false negatives. Additionally, for visualization purposes, we also consider the false positive rate and plot ROC curves to further assess the models' performance.




\section{Experiments: Smart Village Case Study} \label{sec:experiments}

We evaluate our proposed methodology for forecasting repeated visitation in a smart village context. The study area is located in the Alpujarra region of southern Spain and comprises three villages: Pampaneira, Capileira, and Bubion. We represent tourist visitation routes as graphs and perform graph classification to predict whether a tourist will return to the area. 


\subsection{Graphs preprocessing and representation} \label{sec:data_collection_preprocessing}


We collected license plate data using ANPR devices as vehicles entered or exited the area. The IoT scenario, illustrated in Figure \ref{fig:map_as_cameras}, consisted of four Hikvision IP cameras that leverage deep learning-based ANPR technology. Each camera is equipped with vehicle detection sensors, a 2MP resolution, varifocal lenses (2.8–12\,mm), and infrared LEDs with a 50-meter range, ensuring complete coverage of vehicular movements.



The ANPR devices recorded license plate numbers along with timestamps and the movement direction for each vehicle. After data collection, we preprocessed the raw data (see Table \ref{tab:raw_dataset}) by correcting incomplete routes by merging partial capture records based on matching license plates and plausible time gaps, ensuring that each vehicle’s passage through the area was fully reconstructed \cite{duran2024routerecoverer}. In our study, we use four cameras: PAM1 and PAM2 at the entrance and exit of Pampaneira, BUB at the entrance of Bubion, and CAP at the entrance of Capileira.

We recorded data from February 2022 to December 2024; the first six months were used for the experiments, while the remaining 2 years and 5 months were used to label the dataset for validation by determining whether a tourist returned. This dataset covers 467,773 different vehicles and includes entry and exit times, node appearances, travel times between nodes, and the capture direction. We have made this real-world dataset from our smart village area publicly available in \cite{duran_lopez_2025_14993130}.

\begin{table}[htbp]
\centering
\caption{Raw Dataset}
\label{tab:raw_dataset}
\footnotesize
\begin{tabular}{lll}
\toprule
\textbf{Variable} & \textbf{Description} & \textbf{Source} \\ 
\midrule
num\_plate & License plate number & ANPR Cameras \\ 
camera\_ID & Identifier for ANPR camera & ANPR Cameras \\ 
date & Timestamp of appearance & ANPR Cameras \\ 
direction & Direction of vehicle (entry or exit) & ANPR Cameras \\ 
\bottomrule
\end{tabular}%
\end{table}




We preprocess the data to represent our study area as a graph (see Figure \ref{fig:map_as_graphs}). In this representation, each etourist route is represented as a graph $G = (V, E, X_v, X_e)$. 

\begin{itemize}
	\item Each of the four cameras becomes a \textbf{node} $v \in V = \{\text{PAM1, PAM2, BUB, CAP}\}$ and is associated with a feature vector $x_v \in \mathbb{R}^d$, obtained by concatenating a one‐hot encoding of the camera identifier with a directional attribute, indicating whether the vehicle is entering (1) or exiting (0) through that camera. In this way, $X_v = \{\,x_v : v \in V\}$ collects all node feature vectors.
	\item To form \textbf{edges} $(v_i, v_j) \in E$, we connect cameras in the exact sequence that a vehicle is detected; for example, if a vehicle appears first at PAM2 and then at BUB, there is a directed edge from PAM2 to BUB. Each edge carries an attribute $x_{e_{ij}} \in \mathbb{R}^p$ equal to the travel time between those two detections, calculated as the difference between their timestamps. $X_e = \{\,x_{e_{ij}} : (v_i, v_j) \in E\}$ collects all edge feature vectors. Because direction is recorded at each node, all edges are directed.
\end{itemize}

Finally, linking all camera detections builds the graph of the vehicle’s route trip. This fixed graph representation preserves the sequential pattern of each route and serves as the input for our binary graph classification task, where the target variable $Y$ indicates whether the tourist will revisit the area. A tourist is considered a non-repeater if they do not return (only one visit is recorded), and a repeater if they make more than one visit in the future.







\begin{center}
  \includegraphics[width=1\linewidth]{Images/map_Grafos_0.png}
  \captionof{figure}{Study area with the ANPR position}
  \label{fig:map_as_cameras}
\end{center}

\begin{center}
  \includegraphics[width=1\linewidth]{Images/map_Grafos_1.png}
  \captionof{figure}{Study area as a graph}
  \label{fig:map_as_graphs}
\end{center}


\subsection{Graph Neural Network} \label{sec:graph_classification}

Using these graph objects, we perform the graph classification task to predict whether a tourist will revisit the area using solely the graph information. To address this task, we train different GNNs to evaluate their performance on this task, including spectral-based models such as GCN \cite{kipf2016semi}, spatial-based models like GraphSAGE \cite{hamilton2017inductive}, EdgeConv \cite{wang2019dynamic}, and MPNN \cite{gilmer2017neural}, attention-based models such as GAT \cite{velivckovic2017graph}, and recurrent models like GGNN \cite{li2015gated}. In addition, we incorporate their edge-enhanced variants E-GCN \cite{gong2019exploiting}, E-GraphSAGE \cite{lo2022graphsage}, E-GAT \cite{wang2021egat, chen2021edge}, and E-GGNN \cite{zheng2024edge, jiao2021edge} that leverage edge feature information to enrich the representation learning process, allowing us to assess the impact of incorporating edge features on classification performance. 


\subsection{Explainability} \label{sec:explanaibility}


We apply GNNExplainer to our best-performing model to uncover the importance of both node and edge features in driving its predictions. By generating soft masks for nodes and edges, the explainer identifies which features and connections are most influential, allowing us to determine whether incorporating edge features yields a significant benefit. 



\subsection{Neural Architecture Search} \label{sec:nas_section}


For NAS-HPO (see Table \ref{tab:hyperparameters}), we optimize our best performing GNN architecture by tuning both the layer configurations and training hyperparameters, following the settings detailed in Sections \ref{sec:nas_subsec} and \ref{sec:model_training_eval}. We employ Optuna \cite{akiba2019optuna} using the TPESampler along with a Hyperband pruner, running 200 trials. In each trial, we train the model for 30 epochs with a batch size of 64, with early stopping triggering if the validation F1-score does not improve over 5 consecutive epochs. 

Based on the evaluations (see Table \ref{tab:classification_report_smart_poqueira}), the E-GAT model outperformed other GNN models. Motivated by these findings, Table \ref{tab:hyperparameters} presents the configurable architecture layers and parameters that define the search space for NAS-HPO. The works \cite{wang2021egat, chen2021edge} present similar E-GAT architectures, both of which include node and edge projection layers, E-GAT layers, and fusion MLP layers; however, we extend the architecture by incorporating residual connections to improve gradient flow and facilitate training \cite{he2016deep}. While these works use sum as the aggregation function, we also experiment with mean and max. Additionally, although they use LeakyReLU, we further explore GELU and ReLU as activation functions. Moreover, none of the papers mention applying dropout within the MLP or attention layers \cite{srivastava2014dropout}, nor Batch Normalization \cite{ioffe2015batch}, which could be beneficial for model regularization and faster convergence.


\begin{table*}[htbp]
    \centering
    \caption{Configurable Parameters}
    \label{tab:hyperparameters}
    \begin{tabular}{llll}
        \toprule
        \textbf{Parameter} & \textbf{Type} & \textbf{Values} & \textbf{Description} \\
        \midrule
        \multicolumn{4}{l}{\textbf{Architecture Layers}} \\
        \midrule
        Node projection layers   & Integer   & 1, 2, 3         & Number of layers in the node projection MLP. \\
        Edge projection layers   & Integer   & 1, 2, 3         & Number of layers in the edge projection MLP. \\
        E-GAT layers              & Integer   & 1, 2, 3, 4      & Number of iterative E-GAT layers (node and edge updates). \\
        Fusion MLP layers        & Integer   & 1, 2, 3         & Layers in the multi-scale fusion MLP. \\
        Message MLP layers       & Integer   & 0, 1, 2, 3      & Layers in the message MLP. \\
        Residual connections     & Boolean   & True, False     & Use residual links in the MLP blocks. \\
        Aggregation              & String    & mean, sum, max  & Aggregation function used in message passing. \\
        Activation               & String    & ReLU, GELU, LeakyReLU & Activation function used in the MLP layers. \\
        \midrule
        \multicolumn{4}{l}{\textbf{Training Hyperparameters}} \\
        \midrule
        Hidden dim               & Integer   & [32, 256]       & Size of hidden representations. \\
        Dropout rate             & Float     & [0.0, 0.5]      & Dropout probability in MLPs. \\
        BatchNorm                & Boolean   & True, False     & Apply Batch Normalization in MLPs. \\
        Learning rate            & Float     & [$10^{-4}$, $10^{-2}$] & Optimizer learning rate. \\
        Undersampling            & Boolean   & True, False     & Balance the training dataset by undersampling the\\& & &  majority class in a 60-40 proportion. \\
        \bottomrule
    \end{tabular}
\end{table*} 


\subsection{Model Training and Evaluation}

In our approach, we employ the AdamW optimizer \cite{loshchilov2017decoupled} together with the BCEWithLogitsLoss function \cite{paszke2019pytorch} for training our GNN models. To address class imbalance, we compute class weights and incorporate them into the loss function. We use a fixed batch size of 64, and train the models for up to 30 epochs, with early stopping triggered after 5 consecutive epochs without improvement. We tune the remaining training hyperparameters of each GNN over the ranges described in Section \ref{sec:model_training_eval} (listed in Table \ref{tab:hyperparameters}) and select the optimal settings using Optuna. We implement our experiments in Python 3.12.3.

The dataset is split in a stratified manner into 70\% for training, 15\% for validation, and 15\% for testing. The validation set is used to monitor the loss and determine the best threshold that maximizes the F1-score based on the precision-recall curve. After training, we save the best-performing model and evaluate it on the test set using the classification report.



\section{Results and discussion} \label{sec:results_discussion}

Table \ref{tab:classification_report_smart_poqueira} summarizes the classification performance for our two classes, reporting the precision, recall, and F1-score for each class along with the weighted averages across all evaluated models. The experimental results indicate that traditional models based solely on node features and graph structure do not perform well. Models such as GCN, GraphSage, and GGNN perform poorly, with precision, recall, and F1-score barely reaching 0.5 on the weighted average. In contrast, modern models that incorporate edge features, such as E-GCN, E-GraphSAGE, and E-GAT, show improved performance across all classes and metrics. Notably, the E-GAT model achieves the best results across all metrics.


\begin{table}[htbp]
    \centering
    \caption{Classification report of all models.}
    \label{tab:classification_report_smart_poqueira}
        \scalebox{0.7}{\begin{tabular}{l l | c c c | c c c}
            \toprule
            \textbf{Model} & \textbf{Class} & \textbf{Precision} & \textbf{Recall} & \textbf{F1-score} & \multicolumn{3}{c}{\textbf{Weighted Avg}} \\
            \cmidrule(lr){6-8}
            & & & & & \textbf{Precision} & \textbf{Recall} & \textbf{F1-score} \\
            \midrule
            \multirow{2}{*}{GCN} 
                & Non-Repeater    & 0.77 & 0.23 & 0.35 & \multirow{2}{*}{0.63} & \multirow{2}{*}{0.45} & \multirow{2}{*}{0.41} \\
                & Repeater        & 0.37 & 0.87 & 0.52 &  &  &  \\
            \midrule
            \multirow{2}{*}{E-GCN} 
                & Non-Repeater    & 0.78 & 0.74 & 0.76 & \multirow{2}{*}{0.69} & \multirow{2}{*}{0.69} & \multirow{2}{*}{0.69} \\
                & Repeater        & 0.53 & 0.59 & \textbf{0.56} &  &  &  \\
            \midrule
            \multirow{2}{*}{GraphSAGE} 
                & Non-Repeater    & 0.71 & 0.83 & 0.76 & \multirow{2}{*}{0.64} & \multirow{2}{*}{0.66} & \multirow{2}{*}{0.64} \\
                & Repeater        & 0.50 & 0.34 & 0.40 &  &  &  \\
            \midrule
            \multirow{2}{*}{E-GraphSAGE} 
                & Non-Repeater    & 0.77 & 0.78 & 0.77 & \multirow{2}{*}{0.69} & \multirow{2}{*}{0.67} & \multirow{2}{*}{0.68} \\
                & Repeater        & 0.61 & 0.52 & \textbf{0.56} &  &  &  \\ 
            \midrule

            \multirow{2}{*}{GGNN} 
                & Non-Repeater    & 0.78 & 0.33 & 0.46 & \multirow{2}{*}{0.65} & \multirow{2}{*}{0.50} & \multirow{2}{*}{0.49} \\
                & Repeater        & 0.39 & 0.82 & 0.53 &  &  &  \\
            \midrule
            \multirow{2}{*}{E-GGNN} 
                & Non-Repeater    & 0.79 & 0.43 & 0.56 & \multirow{2}{*}{0.55} & \multirow{2}{*}{0.60} & \multirow{2}{*}{0.55} \\
                & Repeater        & 0.41 & 0.78 & 0.54 &  &  &  \\
            \midrule 
            \multirow{2}{*}{EdgeConv} 
                & Non-Repeater    & 0.79 & 0.48 & 0.60 & \multirow{2}{*}{0.67} & \multirow{2}{*}{0.57} & \multirow{2}{*}{0.58} \\
                & Repeater        & 0.43 & 0.75 & 0.54 &  &  &  \\
            \midrule
            \multirow{2}{*}{MPNN} 
                & Non-Repeater    & \textbf{0.80} & 0.56 & 0.66 & \multirow{2}{*}{0.68} & \multirow{2}{*}{0.61} & \multirow{2}{*}{0.63} \\
                & Repeater        & 0.46 & 0.72 & \textbf{0.56} &  &  &  \\
            \midrule
            \multirow{2}{*}{GAT} 
                & Non-Repeater    & 0.75 & 0.11 & 0.20 & \multirow{2}{*}{0.62} & \multirow{2}{*}{0.39} & \multirow{2}{*}{0.30} \\
                & Repeater        & 0.35 & \textbf{0.93} & 0.51 &  &  &  \\
            \midrule
            \multirow{2}{*}{E-GAT} 
                & Non-Repeater    & 0.77 & \textbf{0.78} & \textbf{0.77} & \multirow{2}{*}{\textbf{0.70}} & \multirow{2}{*}{\textbf{0.70}} & \multirow{2}{*}{\textbf{0.70}} \\
                & Repeater        & \textbf{0.56} & 0.55 & 0.55 & & & \\
            \bottomrule
        \end{tabular}}
\end{table}

By applying GNNExplainer to the E-GAT model, we assess the relative importance of edges compared to nodes. Figure \ref{fig:node_and_edge_importance} shows the importance of the time intervals between nodes (derived from the edge mask). These results indicate that, in the best-performing E-GAT model, the time feature is the most important, exhibiting more than twice the importance of the most significant node feature (direction). This finding underscores the advantage of integrating edge features into the attention mechanisms for improved representation.

\begin{figure}[htbp]
    \centering
    \includegraphics[width=0.49\textwidth]{Images/combined_importances.png}
    \caption{Combined importance: nodes with edges}
    \label{fig:node_and_edge_importance}
\end{figure}



Regarding our multivariate NAS-HPO, Figure \ref{fig:pareto_front_trials} shows the NAS process minimizing both the negative F1-score and execution time simultaneously. In this figure, blue markers represent individual trial executions, while red stars indicate the Pareto front, which comprises the best-performing configurations. Additionally, Figure \ref{fig:pareto_front_separate} illustrates the NAS execution across all trials and includes regression lines for both the F1-score and execution time. These lines demonstrate that as the number of trials increases, the validation F1-score improves and the execution time decreases. In our experiment, the first 100 trials are primarily dedicated to exploring the configuration space, while the subsequent trials focus on exploiting the promising regions identified during exploration (see Figure \ref{fig:pareto_front_separate}). Overall, these results confirm that the TPESampler effectively optimizes both model performance and computational efficiency when a higher number of trials are performed.




\begin{center}
  \includegraphics[width=1\linewidth]{Images/pareto_front.png}
  \captionof{figure}{Multivariate NAS-HPO executions. Red stars indicate the Pareto front.}
  \label{fig:pareto_front_trials}
\end{center}



\begin{center}
  \includegraphics[width=1\linewidth]{Images/trial_evolution.png}
  \captionof{figure}{NAS-HPO execution evolution across trials. Regression lines indicate trends in both metrics.}
  \label{fig:pareto_front_separate}
\end{center}

We perform one final execution using the best trial from the Pareto front obtained through applying the multivariate NAS-HPO to the E-GAT model. The corresponding optimal layer configurations and hyperparameters are provided in Table \ref{tab:best_trial_parameters}. Notably, E-GAT + NAS-HPO, the extended architecture, incorporates additional elements that were not present in the original model, such as residual connections, a mean-based aggregation function, GELU activation, dropout in the MLP, and Batch Normalization. These modifications yield a new model that outperforms the base architecture. Table \ref{tab:classification_report_smart_poqueira_egats} presents the precision, recall, and F1-score metrics, along with the execution time in seconds for both E-GAT and E-GAT + NAS-HPO, evaluated on the test set. This table highlights that the extended architecture yields a 5.71\% improvement in evaluation metrics, and the execution time is reduced from 141.76 seconds to 112.17 seconds, representing a 20.87\% decrease. This enhancement is further demonstrated in Figure \ref{fig:all_models_roc}, which displays the ROC curves for the E-GAT and E-GAT + NAS-HPO models compared to the baseline GAT model. These results show that incorporating edge features and optimizing the architecture with NAS-HPO leads to a more effective model that improves the graph classification task.




\begin{table}[htbp]
    \centering
    \caption{Parameters of best NAS-HPO architecture}
    \label{tab:best_trial_parameters}
    \begin{tabular}{llll}
    \toprule
 \multicolumn{2}{c}{\textbf{Architecture Layers}}& \multicolumn{2}{c}{\textbf{Training Hyperparameters}}\\
        \midrule
        \textbf{Parameter} & \textbf{Value}  & \textbf{Parameter}&\textbf{Value}\\
        \midrule
        Node projection layers   & 1  & Hidden dim               &144  \\
        Edge projection layers   & 2  & Dropout rate             &0.5  \\
        E-GAT layers              & 1  & BatchNorm                &True  \\
        Fusion MLP layers        & 2  & Learning rate            &0.0013  \\
        Message MLP layers       & 1  & Undersampling            &False  \\
        Residual connections     & True  & &\\
        Aggregation              & mean  & &\\
        Activation               & GELU  & &\\
        \bottomrule
    \end{tabular}
\end{table}



\begin{center}
  \includegraphics[width=1\linewidth]{Images/all_models_roc.png}
  \captionof{figure}{ROC Curve of base GAT and extended models}
  \label{fig:all_models_roc}
\end{center}          




\begin{table*}[htbp]
    \centering
    \caption{Classification report of E-GAT models with execution times.}
    \label{tab:classification_report_smart_poqueira_egats}
    \resizebox{0.8 \textwidth}{!}{%
    \begin{tabular}{l l | c c c | c c c | c}
        \toprule
        \textbf{Model} & \textbf{Class} & \textbf{Precision} & \textbf{Recall} & \textbf{F1-score} & \multicolumn{3}{c}{\textbf{Weighted Avg}} & \textbf{Execution Time (s)} \\
        \cmidrule(lr){6-8}
         &  &  &  &  & \textbf{Precision} & \textbf{Recall} & \textbf{F1-score} &  \\
        \midrule
        \multirow{2}{*}{E-GAT} 
            & Non-Repeater    & 0.77 & 0.78 & 0.77 & \multirow{2}{*}{0.70} & \multirow{2}{*}{0.70} & \multirow{2}{*}{0.70} & \multirow{2}{*}{141.76} \\
            & Repeater        & 0.56 & \textbf{0.55} & 0.55 &  &  &  &  \\
        \midrule
        \multirow{2}{*}{E-GAT + NAS-HPO} 
            & Non-Repeater    & \textbf{0.78} & \textbf{0.86} & \textbf{0.82} & \multirow{2}{*}{\textbf{0.74}} & \multirow{2}{*}{\textbf{0.74}} & \multirow{2}{*}{\textbf{0.74}} & \multirow{2}{*}{\textbf{112.17}} \\
            & Repeater        & \textbf{0.65} & 0.53 & \textbf{0.58} &  &  &  &  \\
        \bottomrule
    \end{tabular}%
    }
\end{table*}





This study delivers substantial benefits across social, economic, and planning domains. By using deep learning to forecast repeat tourism, it guides the formulation of development policies. Return visits show that protecting heritage and ecosystems boosts repeat tourism \cite{bichler2021designing}, drives higher satisfaction, and economic gains \cite{acharya2023linking}, and lets policymakers plan infrastructure and resources efficiently \cite{alfarhan2022first}.





\section{Framework Transferability}


\begin{table}[htbp] 
\centering
\caption{Ecobici Raw Dataset}
\label{tab:raw_dataset_Ecobici}
\footnotesize
\begin{tabular}{ll}
\toprule
\textbf{Variable}          & \textbf{Description}                                   \\ 
\midrule
user\_id                   & Unique identifier for each user                        \\ 
origin\_station\_id        & Identifier of the station where the trip began          \\ 
destination\_station\_id   & Identifier of the station where the trip ended          \\ 
start\_timestamp           & Date and time when the trip started                     \\ 
end\_timestamp             & Date and time when the trip ended                       \\ 
\bottomrule
\end{tabular}
\end{table}




\begin{figure}[htbp]
  \centering
  \includegraphics[width=0.83\linewidth]{Images/ecobici.png}
  \caption{Geographic location of the six stations selected to simulate a smart village environment.}
  \label{fig:station_selection}
\end{figure}

To study the transferability of our framework, due to the lack of publicly available smart village datasets, we had to artificially create one from a smart city dataset. We selected Ecobici \footnote{\url{https://data.buenosaires.gob.ar/dataset/bicicletas-publicas}}, a public dataset of bicycle trips in Buenos Aires, Argentina \cite{sanchez2018modeling}. Although Ecobici spans all neighborhoods across Buenos Aires, we selected six nearby and isolated stations to emulate a smaller, self-contained network, simulating a smart village. Figure \ref{fig:station_selection} shows the locations of these six stations on a map of Buenos Aires. Ecobici records consecutive station-to-station trips tied to unique user IDs, which allows us to convert the dataset into graphs. Ecobici records every trip taken in the Buenos Aires Public Bike System, including trip identifiers, timestamps, station IDs, and user information. Table \ref{tab:raw_dataset_Ecobici} shows the variables used to construct each user’s trip path. The dataset spans from January to December 2024, and comprises 954 unique users. The first six months form our training period, and the last six months provide labels for churn prediction. Concretely, we build one graph per user per month. For each graph, if the user appears in any subsequent month, we label it as non-churn; if the user does not appear in any subsequent month, we label it as churn. We now follow the methodology in \ref{sec:methodology}, which consists of the following steps:

\begin{itemize}
    \item \textbf{Preprocessing:} We begin by preprocessing the raw data to create graph objects. Bicycle stations become nodes $V$ in the graph. We add an edge $(v_i, v_j) \in E$ between two stations the ``user\_id'' visited, with the edge attribute $x_{e_{ij}} \in \mathbb{R}^p$ representing the difference between their respective timestamps of the stations.  We construct a graph by aggregating all trips the user took during that month.
    
    \item \textbf{Graph Neural Network:}  Using these graph objects, we predict user churn based. We apply different GNN models, both with and without edge attributes, to assess how including those attributes affects performance. Table \ref{tab:classification_report_ecobici} shows the results.
    \item \textbf{Explainability:} To understand which nodes and edges most significantly affect the model's predictions we apply explainability and compute masks for nodes and edges. Figure \ref{fig:node_and_edge_importance_ecobici} helps us identify the global importance score for the nodes and the edge-time feature.
    \item \textbf{Neural Architecture Search:} After selecting the best-performing GNN model, we apply the NAS‐HPO approach. Table \ref{tab:classification_report_ecobici_egats} shows the resulting classification metrics, and Table \ref{tab:hyperparameters_below} lists the optimal network architecture along with its hyperparameters.
\end{itemize}


Table \ref{tab:classification_report_ecobici_egats} shows that GNNs incorporating edge attributes outperform node-only models, with E-GAT achieving the highest metrics. After applying NAS-HPO to the E-GAT base model, we observe an 11.11\% increase in F1-Score and a 54.18\% reduction in execution time. 

The Ecobici dataset demonstrates the transferability of our framework to a churn prediction in a bicycle-sharing service; distinct from our return prediction in a touristic village. It also shows that our methodology highlights how GNNs with edge attributes excel on datasets featuring temporal information, underscoring the importance of incorporating travel‐time data when designing GNNs for mobility analysis. In both cases, E‐GAT outperforms other architectures, and E-GAT with NAS‐HPO further improves its performance.






\begin{table}[htbp]
    \centering
    \caption{Classification report of all models in Ecobici dataset.}
    \label{tab:classification_report_ecobici}
    \scalebox{0.7}{\begin{tabular}{l l | c c c | c c c}
        \toprule
        \textbf{Model} & \textbf{Class} & \textbf{Precision} & \textbf{Recall} & \textbf{F1‐score} & \multicolumn{3}{c}{\textbf{Weighted Avg}} \\
        \cmidrule(lr){6-8}
        & & & & & \textbf{Precision} & \textbf{Recall} & \textbf{F1‐score} \\
        \midrule
        \multirow{2}{*}{GCN} 
            & Non‐Churn    & 0.48 & 0.23 & 0.31 & 0.41 & 0.35 & 0.34 \\
            & Churn        & 0.29 & 0.56 & 0.38 &      &      &      \\

        \midrule
        \multirow{2}{*}{E‐GCN} 
            & Non‐Churn  & 0.70 & 0.43 & 0.53 & 0.59 & 0.52 & 0.52 \\
            & Churn      & 0.40 & 0.68 & 0.50 &      &      &      \\
        \midrule
        \multirow{2}{*}{GraphSAGE} 
        & Non‐Churn & \textbf{0.78} & 0.17 & 0.28 & 0.62 & 0.44 & 0.37 \\
        & Churn     & 0.38 & \textbf{0.91} & 0.54 &      &      &      \\
        \midrule
        \multirow{2}{*}{E‐GraphSAGE} 
            & Non‐Churn  & 0.71 & 0.38 & 0.50 & 0.60 & 0.51 & 0.50 \\
            & Churn      & 0.40 & 0.72 & 0.52 &      &      &      \\
        \midrule
        \multirow{2}{*}{GGNN} 
            & Non‐Churn & 0.72 & 0.26 & 0.38 & 0.60 & 0.46 & 0.43 \\
            & Churn     & 0.39 & 0.82 & 0.52 &      &      &      \\
        \midrule
        \multirow{2}{*}{E‐GGNN} 
            & Non‐Churn & 0.72 & 0.38 & 0.50 & 0.61 & 0.51 & 0.51 \\
            & Churn     & 0.41 & 0.74 & \textbf{0.53} &      &      &      \\
        \midrule
        \multirow{2}{*}{EdgeConv} 
            & Non‐Churn & 0.71 & 0.54 & 0.61 & 0.61 & 0.56 & 0.57 \\
            & Churn     & 0.43 & 0.61 & 0.50 &      &      &      \\

        \midrule
        \multirow{2}{*}{MPNN} 
            & Non‐Churn & 0.68 & 0.17 & 0.28 & 0.57 & 0.42 & 0.36 \\
            & Churn     & 0.37 & 0.86 & 0.52 &      &      &      \\

        \midrule
        \multirow{2}{*}{GAT} 
            & Non‐Churn     & 0.70 & 0.44 & 0.54 & 0.59 & 0.52 & 0.53 \\
            & Churn         & 0.40 & 0.67 & 0.50 &      &      &      \\
        \midrule
        \multirow{2}{*}{E‐GAT} 
            & Non‐Churn    & 0.72 & \textbf{0.68} & \textbf{0.70} & \textbf{0.63} & \textbf{0.62} & \textbf{0.63} \\
            & Churn         & \textbf{0.48} & 0.52 & 0.50 &      &      &      \\
        \bottomrule
    \end{tabular}}
\end{table}


\begin{table*}[htbp]
    \centering
    \caption{Classification report of E-GAT models with execution times in Ecobici dataset.}
    \label{tab:classification_report_ecobici_egats}
    \resizebox{0.8\textwidth}{!}{%
    \begin{tabular}{l l | c c c | c c c | c}
        \toprule
        \textbf{Model} & \textbf{Class} & \textbf{Precision} & \textbf{Recall} & \textbf{F1-score} & \multicolumn{3}{c}{\textbf{Weighted Avg}} & \textbf{Execution Time (s)} \\
        \cmidrule(lr){6-8}
         &  &  &  &  & \textbf{Precision} & \textbf{Recall} & \textbf{F1-score} &  \\
        \midrule
        \multirow{2}{*}{E-GAT} 
            & Non-Churn    & \textbf{0.72} & \textbf{0.68} & \textbf{0.70} & \multirow{2}{*}{0.63} & \multirow{2}{*}{0.62} & \multirow{2}{*}{0.63} & \multirow{2}{*}{26.54} \\
            & Churn        & 0.48 & 0.52 & 0.50 &  &  &  &  \\
        \midrule
        \multirow{2}{*}{E-GAT + NAS-HPO} 
            & Non-Churn    & 0.67 & 0.66 & \textbf{0.70} & \multirow{2}{*}{\textbf{0.67}} & \multirow{2}{*}{\textbf{0.66}} & \multirow{2}{*}{\textbf{0.70}} & \multirow{2}{*}{\textbf{12.16}} \\
            & Churn        & \textbf{0.66} & \textbf{0.65} & \textbf{0.69} &  &  &  &  \\
        \bottomrule
    \end{tabular}%
    }
\end{table*}



\begin{figure}[htbp]
    \centering
    \includegraphics[width=0.49\textwidth]{Images/combined_importances_transfer.png}
    \caption{Combined importance: nodes with edges}

    \label{fig:node_and_edge_importance_ecobici}
\end{figure}


\begin{table}[htbp]
    \centering
    \caption{Parameters of the best NAS-HPO architecture}
    \label{tab:hyperparameters_below}
    \begin{tabular}{llll}
    \toprule
 \multicolumn{2}{c}{\textbf{Architecture Layers}}& \multicolumn{2}{c}{\textbf{Training Hyperparameters}}\\
        \midrule
        \textbf{Parameter} & \textbf{Value}  & \textbf{Parameter}&\textbf{Value}\\
        \midrule
        Node projection layers   & 1  & Hidden dim               &144  \\
        Edge projection layers   & 1  & Dropout rate             &0.5  \\
        E-GAT layers              & 1  & BatchNorm                &True  \\
        Fusion MLP layers        & 1  & Learning rate            &0.029  \\
        Message MLP layers       & 1  & Undersampling            &False  \\
        Residual connections     & True  & &\\
        Aggregation              & max  & &\\
        Activation               & GELU  & &\\
        \bottomrule
    \end{tabular}
\end{table}




\section{Conclusion} \label{sec:conclusion}

This work presents a framework for modeling and predicting tourist routes in smart villages through the application of GNNs combined with explainability techniques and a multivariate NAS-HPO to forecast repeat tourist visits in a smart village context. By transforming vehicle detection data from IoT ANPR devices into graph representations, we capture detailed spatiotemporal patterns inherent in vehicular movements. By integrating edge features into these representations, we extend GNNs to capture intricate spatial and temporal patterns. 

The results show significant improvements in model performance when incorporating edge features. This was confirmed by the use of explainability techniques, specifically GNNExplainer, which provided clear insights into feature importance and validated the benefits of integrating edge attributes to enhance model performance. When applying multivariate NAS-HPO, we extended the E-GAT model by adding new parameters, achieving a 5.71\% improvement in F1-score and a 20.87\% decrease in execution time for the graph classification task, marking a valuable contribution to smart village management and tourism analytics. Finally, we studied the framework’s transferability, demonstrating that our proposed methodology can be applied in other contexts.



% if have a single appendix:
%\appendix[Proof of the Zonklar Equations]
% or
%\appendix  % for no appendix heading
% do not use \section anymore after \appendix, only \section*
% is possibly needed

% use appendices with more than one appendix
% then use \section to start each appendix
% you must declare a \section before using any
% \subsection or using \label (\appendices by itself
% starts a section numbered zero.)
%


%\appendices
%\section{Proof of the First Zonklar Equation}
%Appendix one text goes here.

% you can choose not to have a title for an appendix
% if you want by leaving the argument blank
%\section{hello}
%Appendix two text goes here.


% use section* for acknowledgment

\section*{Acknowledgment}

This publication is part of the projects C-SEJ-128-UGR23 and PID2023-149185OB-I00 funded by MICIU/AEI/10.13039/501100011033 and by ERDF/EU.


% Can use something like this to put references on a page
% by themselves when using endfloat and the captionsoff option.
\ifCLASSOPTIONcaptionsoff
  \newpage
\fi



% trigger a \newpage just before the given reference
% number - used to balance the columns on the last page
% adjust value as needed - may need to be readjusted if
% the document is modified later
%\IEEEtriggeratref{8}
% The "triggered" command can be changed if desired:
%\IEEEtriggercmd{\enlargethispage{-5in}}

% references section

% can use a bibliography generated by BibTeX as a .bbl file
% BibTeX documentation can be easily obtained at:
% http://mirror.ctan.org/biblio/bibtex/contrib/doc/
% The IEEEtran BibTeX style support page is at:
% http://www.michaelshell.org/tex/ieeetran/bibtex/
\bibliographystyle{IEEEtran}
% argument is your BibTeX string definitions and bibliography database(s)
\bibliography{IEEEabrv, cas-refs}

%
% <OR> manually copy in the resultant .bbl file
% set second argument of \begin to the number of references
% (used to reserve space for the reference number labels box)

% biography section
% 
% If you have an EPS/PDF photo (graphicx package needed) extra braces are
% needed around the contents of the optional argument to biography to prevent
% the LaTeX parser from getting confused when it sees the complicated
% \includegraphics command within an optional argument. (You could create
% your own custom macro containing the \includegraphics command to make things
% simpler here.)
%\begin{IEEEbiography}[{\includegraphics[width=1in,height=1.25in,clip,keepaspectratio]{mshell}}]{Michael Shell}
% or if you just want to reserve a space for a photo:

%\begin{IEEEbiography}{Michael Shell}
%Biography text here.
%\end{IEEEbiography}


% You can push biographies down or up by placing
% a \vfill before or after them. The appropriate
% use of \vfill depends on what kind of text is
% on the last page and whether or not the columns
% are being equalized.

%\vfill

% Can be used to pull up biographies so that the bottom of the last one
% is flush with the other column.
%\enlargethispage{-5in}



% that's all folks
\end{document}